\section{Industrial \& Commercial Applications of Embedded Linux}

While hobbyists are familiar with the Raspberry Pi, it is crucial to differentiate between educational prototyping hardware and industrial-grade embedded systems.

\subsection{Where Embedded Linux is Deployed}
\begin{itemize}[leftmargin=*]
    \item \textbf{Industrial Networking \& Infrastructure:} Commercial routers, IoT edge gateways, firewalls, and remote telemetry units (RTUs).
    \item \textbf{Aerospace \& Defense:} Autonomous vehicle controllers, mission systems, and spaceflight computers (e.g., SpaceX Falcon 9 rocket flight computers run custom Linux kernels).
    \item \textbf{Robotics \& Automation:} ROS (Robot Operating System) nodes, automated guided vehicles (AGVs), and industrial manipulator arms.
    \item \textbf{Consumer Electronics \& Point-of-Sale (POS):} Smart appliances, interactive kiosks, digital signage, and medical monitoring equipment.
\end{itemize}

\begin{warningbox}{Why the Raspberry Pi is NOT an Industrial Solution}
The Raspberry Pi is an outstanding educational tool and desktop replacement, but it is rarely used in commercial products due to:
\begin{enumerate}[leftmargin=*]
    \item \textbf{Software Bloat:} General-purpose Raspberry Pi OS images contain GUI desktop suites, background daemons, and packages that cause slow boot times and high power draw.
    \item \textbf{Proprietary GPU Boot Architecture:} Booting relies on closed-source VideoCore GPU binary blobs rather than open standard U-Boot flows.
    \item \textbf{Hardware Availability \& Cost:} Lack of long-term industrial supply chain guarantees (10+ year availability required for industrial boards).
\end{enumerate}
\textbf{Production Approach:} Engineers design around System-on-Modules (SOMs) or reference SBC architectures (e.g., STM32MP1, NXP i.MX8, Microchip SAMA5) running a stripped-down, custom-built Linux distribution.
\end{warningbox}
